\documentclass[11pt,a4paper]{article}
\usepackage[top=1.6cm,bottom=1.4cm,left=1.3cm,right=1.3cm,headsep=8pt,footskip=16pt]{geometry}
\usepackage{enumitem}
\usepackage{amssymb}
\usepackage{array}
\usepackage[table]{xcolor}
\usepackage{fancyhdr}
\usepackage{multicol}
\usepackage{polyglossia}
\setdefaultlanguage[calendar=gregorian, numerals=maghrib]{arabic}
\setotherlanguage{english}
\IfFontExistsTF{Amiri}{%
  \newfontfamily\arabicfont[Script=Arabic, Scale=1.08]{Amiri}
}{%
  \newfontfamily\arabicfont[Script=Arabic, Scale=1.08, Path=../fonts/, Extension=.ttf,
    UprightFont=*-Regular, BoldFont=*-Bold, ItalicFont=*-Italic,
    BoldItalicFont=*-BoldItalic]{Amiri}
}
\IfFontExistsTF{AlBattar}{%
  \newfontfamily\khatt[Script=Arabic, Scale=1.15]{AlBattar}
}{%
  \newfontfamily\khatt[Script=Arabic, Scale=1.15, Path=../fonts/,
    BoldFont=ae_AlBattar.ttf]{ae_AlBattar.ttf}
}
\renewcommand{\headrulewidth}{0.8pt}
\renewcommand{\footrulewidth}{0.8pt}
\fancyhf{}
\rhead{\khatt المدرسة العليا للأساتذة ببوزريعة}
\lhead{\khatt ملحقة بومرداس}
\lfoot{\khatt الأستاذ: د. بن عبد الرزاق يوسف}
\cfoot{\thepage/2}
\pagestyle{fancy}
\parindent 0pt
\setlength{\columnsep}{0.6cm}
\setlength{\columnseprule}{0.4pt}
\setlist[enumerate,1]{label=\textbf{\arabic*.}, leftmargin=*, itemsep=5pt, topsep=2pt}
\begin{document}

\begin{center}
{\Large\khatt\textbf{الفرض السادس في الذكاء الاصطناعي}}\\[2pt]
{\normalsize أسئلة متعددة الاختيارات ــ الفصول من 1 إلى 6 ــ 15 سؤالًا}
\end{center}
\vspace{-4pt}\rule{\textwidth}{0.6pt}\\
{\small\textbf{التعليمة:} لكل سؤال إجابة صحيحة واحدة فقط من بين خمسة اختيارات.
اكتب \textbf{حرف} الإجابة الصحيحة في جدول الإجابات أسفل الصفحة الثانية.
لكل سؤال نقطة واحدة (المجموع 15 نقطة).}
\vspace{-2pt}
\begin{multicols}{2}

\noindent\begin{minipage}{\linewidth}
\hangindent=0.7cm\hangafter=1\makebox[0.7cm][r]{\textbf{1.}}\textbf{كم عدد مراحل تعلّم الآلة حسب الدرس؟}
\begin{itemize}[leftmargin=0.95cm, labelsep=4pt, itemsep=0pt, topsep=2pt, parsep=0pt]
  \item[\textbf{أ)}] سبع مراحل
  \item[\textbf{ب)}] ست مراحل
  \item[\textbf{ج)}] ثلاث مراحل
  \item[\textbf{د)}] خمس مراحل
  \item[\textbf{هـ)}] أربع مراحل
\end{itemize}
\end{minipage}\par\vspace{7pt plus 4pt}
\noindent\begin{minipage}{\linewidth}
\hangindent=0.7cm\hangafter=1\makebox[0.7cm][r]{\textbf{2.}}\textbf{ما الإفصاح؟}
\begin{itemize}[leftmargin=0.95cm, labelsep=4pt, itemsep=0pt, topsep=2pt, parsep=0pt]
  \item[\textbf{أ)}] التصريح باستعمال أداة ذكاء اصطناعي
  \item[\textbf{ب)}] نسخ النص الذي أنتجته الأداة كما هو
  \item[\textbf{ج)}] ترجمة النص المكتوب إلى لغة أخرى
  \item[\textbf{د)}] حذف كل المصادر من آخر البحث
  \item[\textbf{هـ)}] إخفاء استعمال الأداة عن الأستاذ دائمًا
\end{itemize}
\end{minipage}\par\vspace{7pt plus 4pt}
\noindent\begin{minipage}{\linewidth}
\hangindent=0.7cm\hangafter=1\makebox[0.7cm][r]{\textbf{3.}}\textbf{تصوير صفحة مخطوطة وطلب كتابتها نصًا رقميًا هو تحويل:}
\begin{itemize}[leftmargin=0.95cm, labelsep=4pt, itemsep=0pt, topsep=2pt, parsep=0pt]
  \item[\textbf{أ)}] نص $\leftarrow$ مخطط
  \item[\textbf{ب)}] نص $\leftarrow$ صوت
  \item[\textbf{ج)}] نص $\leftarrow$ صورة
  \item[\textbf{د)}] صورة $\leftarrow$ نص
  \item[\textbf{هـ)}] صوت $\leftarrow$ فيديو
\end{itemize}
\end{minipage}\par\vspace{7pt plus 4pt}
\noindent\begin{minipage}{\linewidth}
\hangindent=0.7cm\hangafter=1\makebox[0.7cm][r]{\textbf{4.}}\textbf{ما الترتيب الصحيح لعمل الأداة التوليدية؟}
\begin{itemize}[leftmargin=0.95cm, labelsep=4pt, itemsep=0pt, topsep=2pt, parsep=0pt]
  \item[\textbf{أ)}] محتوى جديد $\leftarrow$ الأداة تحلل الطلب $\leftarrow$ المستخدم يكتب طلبًا
  \item[\textbf{ب)}] الأداة تحلل الطلب $\leftarrow$ محتوى جديد $\leftarrow$ المستخدم يكتب طلبًا
  \item[\textbf{ج)}] المستخدم يكتب طلبًا $\leftarrow$ الأداة تحلل الطلب $\leftarrow$ محتوى جديد
  \item[\textbf{د)}] الأداة تحلل الطلب $\leftarrow$ المستخدم يكتب طلبًا $\leftarrow$ محتوى جديد
  \item[\textbf{هـ)}] المستخدم يكتب طلبًا $\leftarrow$ محتوى جديد $\leftarrow$ الأداة تحلل الطلب
\end{itemize}
\end{minipage}\par\vspace{7pt plus 4pt}
\noindent\begin{minipage}{\linewidth}
\hangindent=0.7cm\hangafter=1\makebox[0.7cm][r]{\textbf{5.}}\textbf{لماذا نختبر النموذج ببيانات جديدة لم يرها؟}
\begin{itemize}[leftmargin=0.95cm, labelsep=4pt, itemsep=0pt, topsep=2pt, parsep=0pt]
  \item[\textbf{أ)}] لزيادة عدد الأمثلة في التدريب
  \item[\textbf{ب)}] لنعرف هل فهم حقًا، لا هل حفظ
  \item[\textbf{ج)}] لأن بيانات التدريب خاطئة دائمًا
  \item[\textbf{د)}] لأن بيانات التدريب تُحذف بعده
  \item[\textbf{هـ)}] لتسريع مرحلة التدريب الطويلة
\end{itemize}
\end{minipage}\par\vspace{7pt plus 4pt}
\noindent\begin{minipage}{\linewidth}
\hangindent=0.7cm\hangafter=1\makebox[0.7cm][r]{\textbf{6.}}\textbf{ما السؤال الذي يجيب عنه عنصر «المهمة»؟}
\begin{itemize}[leftmargin=0.95cm, labelsep=4pt, itemsep=0pt, topsep=2pt, parsep=0pt]
  \item[\textbf{أ)}] ما المعلومات التي تساعدها على الفهم؟
  \item[\textbf{ب)}] من تكون الأداة وما خبرتها؟
  \item[\textbf{ج)}] ما الشروط التي يجب احترامها؟
  \item[\textbf{د)}] كيف نريد أن يكون شكل الإجابة؟
  \item[\textbf{هـ)}] ماذا نريد أن تفعل الأداة؟
\end{itemize}
\end{minipage}\par\vspace{7pt plus 4pt}
\noindent\begin{minipage}{\linewidth}
\hangindent=0.7cm\hangafter=1\makebox[0.7cm][r]{\textbf{7.}}\textbf{لماذا يهم سلوك أستاذ المستقبل مع أدوات الذكاء الاصطناعي؟}
\begin{itemize}[leftmargin=0.95cm, labelsep=4pt, itemsep=0pt, topsep=2pt, parsep=0pt]
  \item[\textbf{أ)}] لأن القانون يمنعه من استعمالها
  \item[\textbf{ب)}] لأن التلاميذ لا يستعملون الأدوات
  \item[\textbf{ج)}] لأنه قدوة لتلاميذه في سلوكه
  \item[\textbf{د)}] لأن الأداة تراقب سلوك الأستاذ
  \item[\textbf{هـ)}] لأنه يدفع ثمن الأداة من ماله
\end{itemize}
\end{minipage}\par\vspace{7pt plus 4pt}
\noindent\begin{minipage}{\linewidth}
\hangindent=0.7cm\hangafter=1\makebox[0.7cm][r]{\textbf{8.}}\textbf{في خطوات التلخيص، ما الخطوة التي تأتي بعد «حدّد الأفكار الأساسية»؟}
\begin{itemize}[leftmargin=0.95cm, labelsep=4pt, itemsep=0pt, topsep=2pt, parsep=0pt]
  \item[\textbf{أ)}] اذكر المصدر الأصلي
  \item[\textbf{ب)}] اكتب بأسلوبك الخاص
  \item[\textbf{ج)}] قيّم المصادر المختلفة
  \item[\textbf{د)}] اقرأ النص كاملًا
  \item[\textbf{هـ)}] احذف التكرار والأمثلة
\end{itemize}
\end{minipage}\par\vspace{7pt plus 4pt}
\noindent\begin{minipage}{\linewidth}
\hangindent=0.7cm\hangafter=1\makebox[0.7cm][r]{\textbf{9.}}\textbf{ما الترتيب الصحيح لخطوات فهم صورة بالأداة؟}
\begin{itemize}[leftmargin=0.95cm, labelsep=4pt, itemsep=0pt, topsep=2pt, parsep=0pt]
  \item[\textbf{أ)}] ارفعها مع أمر دقيق $\leftarrow$ اختر صورة واضحة $\leftarrow$ قارن بالصورة $\leftarrow$ اقرأ الشرح
  \item[\textbf{ب)}] قارن بالصورة $\leftarrow$ اقرأ الشرح $\leftarrow$ اختر صورة واضحة $\leftarrow$ ارفعها مع أمر دقيق
  \item[\textbf{ج)}] اختر صورة واضحة $\leftarrow$ ارفعها مع أمر دقيق $\leftarrow$ اقرأ الشرح $\leftarrow$ قارن بالصورة
  \item[\textbf{د)}] اقرأ الشرح $\leftarrow$ اختر صورة واضحة $\leftarrow$ ارفعها مع أمر دقيق $\leftarrow$ قارن بالصورة
  \item[\textbf{هـ)}] اختر صورة واضحة $\leftarrow$ قارن بالصورة $\leftarrow$ ارفعها مع أمر دقيق $\leftarrow$ اقرأ الشرح
\end{itemize}
\end{minipage}\par\vspace{7pt plus 4pt}
\noindent\begin{minipage}{\linewidth}
\hangindent=0.7cm\hangafter=1\makebox[0.7cm][r]{\textbf{10.}}\textbf{حسب مخطط الدرس، المبتدأ:}
\begin{itemize}[leftmargin=0.95cm, labelsep=4pt, itemsep=0pt, topsep=2pt, parsep=0pt]
  \item[\textbf{أ)}] اسم منصوب
  \item[\textbf{ب)}] جملة فعلية
  \item[\textbf{ج)}] اسم مجرور
  \item[\textbf{د)}] فعل ماضٍ
  \item[\textbf{هـ)}] اسم مرفوع
\end{itemize}
\end{minipage}\par\vspace{7pt plus 4pt}
\noindent\begin{minipage}{\linewidth}
\hangindent=0.7cm\hangafter=1\makebox[0.7cm][r]{\textbf{11.}}\textbf{كتب أستاذ طلبًا دقيقًا فحصل على إعراب مقنع. ماذا يفعل قبل استعماله في القسم؟}
\begin{itemize}[leftmargin=0.95cm, labelsep=4pt, itemsep=0pt, topsep=2pt, parsep=0pt]
  \item[\textbf{أ)}] يتحقق منه في كتاب نحو معتمد، ثم يقرّر
  \item[\textbf{ب)}] يستعمله مباشرة لأن طلبه كان دقيقًا جدًا
  \item[\textbf{ج)}] يعيد الطلب، فإن تكرر الإعراب فهو صحيح
  \item[\textbf{د)}] يستعمله إذا كان طويلًا ومفصلًا ومنظمًا
  \item[\textbf{هـ)}] يقارنه بأداة ثانية ويعتمد ما اتفقتا عليه
\end{itemize}
\end{minipage}\par\vspace{7pt plus 4pt}
\noindent\begin{minipage}{\linewidth}
\hangindent=0.7cm\hangafter=1\makebox[0.7cm][r]{\textbf{12.}}\textbf{أي سلسلة تصف كيف ينشأ الانحياز الخوارزمي؟}
\begin{itemize}[leftmargin=0.95cm, labelsep=4pt, itemsep=0pt, topsep=2pt, parsep=0pt]
  \item[\textbf{أ)}] بيانات ناقصة $\leftarrow$ نموذج تعلّم صورة ناقصة $\leftarrow$ نتائج منحازة
  \item[\textbf{ب)}] بيانات متنوعة $\leftarrow$ نموذج تعلّم صورة كاملة $\leftarrow$ نتائج منحازة
  \item[\textbf{ج)}] نموذج تعلّم صورة كاملة $\leftarrow$ بيانات ناقصة $\leftarrow$ نتائج عادلة
  \item[\textbf{د)}] نتائج منحازة $\leftarrow$ بيانات ناقصة $\leftarrow$ نموذج تعلّم صورة ناقصة
  \item[\textbf{هـ)}] بيانات متوازنة $\leftarrow$ نموذج تعلّم صورة ناقصة $\leftarrow$ نتائج عادلة
\end{itemize}
\end{minipage}\par\vspace{7pt plus 4pt}
\noindent\begin{minipage}{\linewidth}
\hangindent=0.7cm\hangafter=1\makebox[0.7cm][r]{\textbf{13.}}\textbf{في الأمر: «تصرّف كمساعد أكاديمي. أنشئ خطة مراجعة لمادة النحو مدتها 7 أيام، بمعدل ساعتين يوميًا. نظّم الخطة في جدول»، عبارة «مدتها 7 أيام، بمعدل ساعتين يوميًا» تمثّل:}
\begin{itemize}[leftmargin=0.95cm, labelsep=4pt, itemsep=0pt, topsep=2pt, parsep=0pt]
  \item[\textbf{أ)}] المهمة
  \item[\textbf{ب)}] القيود
  \item[\textbf{ج)}] التنسيق
  \item[\textbf{د)}] الدور
  \item[\textbf{هـ)}] السياق
\end{itemize}
\end{minipage}\par\vspace{7pt plus 4pt}
\noindent\begin{minipage}{\linewidth}
\hangindent=0.7cm\hangafter=1\makebox[0.7cm][r]{\textbf{14.}}\textbf{رتّب الاستعمالات من الأكثر قبولًا إلى المرفوض:}
\begin{itemize}[leftmargin=0.95cm, labelsep=4pt, itemsep=0pt, topsep=2pt, parsep=0pt]
  \item[\textbf{أ)}] أنشئ مسودة وأراجعها وأصرّح $\leftarrow$ أنسخ إجابة كاملة باسمي $\leftarrow$ أفهم مفهومًا صعبًا
  \item[\textbf{ب)}] أنسخ إجابة كاملة باسمي $\leftarrow$ أفهم مفهومًا صعبًا $\leftarrow$ أنشئ مسودة وأراجعها وأصرّح
  \item[\textbf{ج)}] أفهم مفهومًا صعبًا $\leftarrow$ أنسخ إجابة كاملة باسمي $\leftarrow$ أنشئ مسودة وأراجعها وأصرّح
  \item[\textbf{د)}] أفهم مفهومًا صعبًا $\leftarrow$ أنشئ مسودة وأراجعها وأصرّح $\leftarrow$ أنسخ إجابة كاملة باسمي
  \item[\textbf{هـ)}] أنسخ إجابة كاملة باسمي $\leftarrow$ أنشئ مسودة وأراجعها وأصرّح $\leftarrow$ أفهم مفهومًا صعبًا
\end{itemize}
\end{minipage}\par\vspace{7pt plus 4pt}
\noindent\begin{minipage}{\linewidth}
\hangindent=0.7cm\hangafter=1\makebox[0.7cm][r]{\textbf{15.}}\textbf{مدونة شخصية لأستاذ مختص في النحو يذكر فيها مراجعه. كيف نتعامل معها؟}
\begin{itemize}[leftmargin=0.95cm, labelsep=4pt, itemsep=0pt, topsep=2pt, parsep=0pt]
  \item[\textbf{أ)}] نرفضها تمامًا لأن كل المدونات خاطئة وغير مفيدة
  \item[\textbf{ب)}] نقيّمها بالمعايير، ونتحقق منها لأنها أقل ثقة من المحكّم
  \item[\textbf{ج)}] نعدّها مقالًا علميًا محكّمًا لأنها تذكر مراجعها
  \item[\textbf{د)}] لا نطبّق عليها أي معيار لأنها كُتبت من أستاذ
  \item[\textbf{هـ)}] نعدّها أوثق من الكتب المحققة لأن كاتبها مختص
\end{itemize}
\end{minipage}\par\vspace{7pt plus 4pt}
\end{multicols}
\vfill
\begin{center}
{\khatt\textbf{جدول الإجابات}}\\[3pt]
\renewcommand{\arraystretch}{1.6}
\begin{tabular}{|>{\centering\arraybackslash}m{1.9cm}|*{15}{>{\centering\arraybackslash}m{0.62cm}|}}
\hline
\textbf{رقم السؤال} & \textbf{1} & \textbf{2} & \textbf{3} & \textbf{4} & \textbf{5} & \textbf{6} & \textbf{7} & \textbf{8} & \textbf{9} & \textbf{10} & \textbf{11} & \textbf{12} & \textbf{13} & \textbf{14} & \textbf{15} \\ \hline
\textbf{الإجابة} &  &  &  &  &  &  &  &  &  &  &  &  &  &  &  \\[8pt] \hline
\end{tabular}

\vspace{8pt}
\renewcommand{\arraystretch}{2.2}
\begin{tabular}{|>{\centering\arraybackslash}m{7.5cm}|>{\centering\arraybackslash}m{3.2cm}|>{\centering\arraybackslash}m{3.2cm}|}
\hline
\textbf{الاسم واللقب} & \textbf{الفوج} & \textbf{العلامة} \\ \hline
 & & \hfill /15 \\ \hline
\end{tabular}
\end{center}
\end{document}
